\documentclass[11pt]{amsart}

\usepackage{amsmath,amssymb,amsthm}
\usepackage{mathtools}
\usepackage[T1]{fontenc}
\usepackage{stmaryrd} % \llbracket \rrbracket
\usepackage{booktabs}
\usepackage[hidelinks]{hyperref}
\usepackage{enumitem}

% ---- theorem environments ----
\theoremstyle{plain}
\newtheorem{theorem}{Theorem}[section]
\newtheorem{proposition}[theorem]{Proposition}
\newtheorem{lemma}[theorem]{Lemma}
\newtheorem{corollary}[theorem]{Corollary}
\theoremstyle{definition}
\newtheorem{definition}[theorem]{Definition}
\newtheorem{example}[theorem]{Example}
\newtheorem{construction}[theorem]{Construction}
\theoremstyle{remark}
\newtheorem{remark}[theorem]{Remark}
\newtheorem{question}[theorem]{Question}

% ---- notation ----
\newcommand{\Set}{\mathbf{Set}}
\newcommand{\Cont}{\mathbf{Cont}}
\newcommand{\Cat}{\mathbf{Cat}}
\newcommand{\Vect}{\mathbf{Vec}}
\newcommand{\Rel}{\mathbf{Rel}}
\newcommand{\Par}{\mathbf{Par}}
\newcommand{\Stoch}{\mathbf{Stoch}}
\newcommand{\Fam}{\mathrm{Fam}}
\newcommand{\Kl}{\mathrm{Kl}}
\newcommand{\Nat}{\mathrm{Nat}}
\newcommand{\Ran}{\mathrm{Ran}}
\newcommand{\op}{\mathrm{op}}
\newcommand{\ext}[1]{\llbracket #1 \rrbracket}
\newcommand{\extM}[1]{\llbracket #1 \rrbracket_{M}}
\newcommand{\extC}[1]{\llbracket #1 \rrbracket_{\C}}
\newcommand{\extKl}[1]{\llbracket #1 \rrbracket^{\Kl}}
\newcommand{\id}{\mathrm{id}}
\newcommand{\Id}{\mathrm{Id}}
\newcommand{\tri}{\triangleleft}
\newcommand{\triM}{\triangleleft_{M}}
\newcommand{\ceil}[1]{\lceil #1 \rceil}
\newcommand{\y}{\mathsf{y}}
\newcommand{\Writer}{\mathsf{Writer}}
\newcommand{\Maybe}{\mathsf{Maybe}}
\newcommand{\Reader}{\mathsf{Reader}}
\newcommand{\Dist}{\mathsf{Dist}}
\newcommand{\List}{\mathsf{List}}
\newcommand{\PP}{\mathcal{P}}
\newcommand{\C}{\mathcal{C}}
\newcommand{\D}{\mathcal{D}}
\newcommand{\CL}{\mathcal{CL}}
\newcommand{\UN}{\mathcal{UN}}
\DeclareMathOperator{\im}{im}

\begin{document}

\title[Change of base for containers]
{Changing the base of a container: \\ six constructions, two invariants, one collapse}

\author{MacBeth}
\address{Kodamai}
\email{}
\date{15 September 2026}

\begin{abstract}
The category of containers is $\Cont = \Fam(\Set^{\op})$; replacing $\Set$ by a category $\C$ along a
functor $F\colon \Set \to \C$ yields $\Fam(\C^{\op})$, whose objects are containers whose positions
are read out in $\C$. This change of base is where effects and directedness enter the theory of
polynomial functors: $\C = \Kl(M)$ gives effectful ($M$-)containers, $\C = \Vect$ gives linear
containers, $\C = \Cat$ recovers directed containers. We survey six ways monad and comonad structure
is transported through this construction, and then isolate the invariant that decides when the
resulting extension into a category of endofunctors is full and faithful. There are three distinct
extensions in play, with three different notions of naturality; keeping them apart resolves an
apparent paradox (the reader monad has a cartesian closed Kleisli category yet its extension is never
full). Our main result is a sharp collapse: over a commutative-monad base, the enriched container
extension is full and faithful for the identity monad and for no other effect. Closedness of the base
makes the enriched extension \emph{exist}; connectedness of the monoidal unit makes it \emph{full};
these axes are independent, and over $\Kl(M)$ they meet only at $M = \Id$.
\end{abstract}

\maketitle

\begin{center}\small
\fbox{\begin{minipage}{0.92\textwidth}\centering
\emph{Author:} MacBeth (Kodamai).\qquad\emph{Date:} 15 September 2026.\\[2pt]
\emph{Corresponds to} \texttt{scotmacbeth/work-in-progress} \emph{content-commit} \texttt{21aef97}.\\[2pt]
\textbf{Draft --- work in progress; Kodamai grant deliverable.} A report on the monad/comonad
constructions on containers, organised around change of base and enrichment; the new contribution is
the enriched-extension collapse of \S\ref{sec:enrichment} (Theorem~\ref{thm:collapse}). The
per-instance results of \S\ref{sec:census} are proved in the cited companion papers and are not
reproduced.
\end{minipage}}
\end{center}

%============================================================
\section{Introduction}
%============================================================

A \emph{container} $(S,P)$ is a set $S$ of shapes together with, for each shape $s\in S$, a set $P_s$
of positions; it presents the polynomial endofunctor
\[
  \ext{S,P}(X) \;=\; \sum_{s\in S} X^{P_s} \;=\; \sum_{s\in S}\Set(P_s, X),
\]
the interpretation of a strictly-positive datatype whose values are ``a shape together with a payload
at each position'' \cite{AAG}. The container category $\Cont$ is equivalent to $\Fam(\Set^{\op})$: an
object is a family of sets indexed by a set, a morphism $(S,P)\to(T,Q)$ is a shape map $S\to T$
together with, contravariantly, a position map back \cite{NiuSpivak}. The extension
$\ext{-}\colon \Cont \to [\Set,\Set]$ is full and faithful onto the polynomial functors: containers
\emph{are} polynomial functors, presented.

\subsection*{Changing the base}
The formula $\ext{S,P}(X)=\sum_s \Set(P_s,X)$ has one degree of freedom that the datatype reading
hides: the hom-set $\Set(P_s,X)$. What if positions are read out not in $\Set$ but in some other
category $\C$ --- so that observing a position may fail, branch, consult a context, cost a resource,
or return a distribution? Formally, fix a functor $F\colon \Set \to \C$ and consider
\[
  \Fam(\C^{\op}),\qquad \extC{S,P}(X) \;=\; \sum_{s\in S}\C(P_s, X)\quad(X\in\C).
\]
This \emph{change of base} is not cosmetic. It is the single construction under which the main
variations of the polynomial-functor calculus are instances of one theory:
\begin{itemize}[leftmargin=1.6em]
  \item $\C = \Set$: ordinary containers.
  \item $\C = \Kl(M)$ for a monad $M$ on $\Set$: \emph{$M$-containers}, whose positions are read in
    the Kleisli category, so that $\extM{S,P}(X)=\sum_s \Set(P_s, MX)$ is the effectful reading.
  \item $\C = \Vect$: linear containers / the $\Vect$-valued theory.
  \item $\C = \Cat$: directed containers, recovering categories \cite{AhmanUustalu}.
\end{itemize}
The point Neil Ghani has pressed is that change of base is the \emph{frame}, and the per-instance
results are details of transport: a functor $F\colon\Set\to\C$ that carries structure, not merely a
relabelling of the base. This report is organised around that frame.

\subsection*{Two questions per construction}
For each way of putting monad or comonad structure on containers, two questions recur:
\begin{enumerate}[label=(\alph*),leftmargin=2.2em]
  \item \textbf{Availability.} Does the construction exist over this base at all?
  \item \textbf{Faithfulness.} Is the extension into endofunctors full and faithful --- i.e. does the
    combinatorial category of containers embed into the semantic category of functors, so that
    ``container morphism'' and ``natural transformation'' coincide?
\end{enumerate}
Full-and-faithfulness is a \emph{diagnostic}, not the headline: it is the sharpest available test of
whether the syntactic category $\Fam(\C^{\op})$ faithfully models its intended semantics. We lead with
the structure of the container category over each base (\S\ref{sec:census}), and then use the
diagnostic to expose what actually controls it (\S\ref{sec:enrichment}).

\subsection*{Results}
The survey (\S\ref{sec:census}) collects the transport constructions Ghani enumerated, each proved
elsewhere in this programme: the position-op monad-to-comonad transfer; the $\forall/\exists$
predicate liftings; the identification of $M$-containers with $\Fam(\Kl(M)^{\op})$ and its codensity
criterion for fullness; algebraic effects as polynomial functors, with the free monad on a container
as their carrier; and the composition-versus-representability dividing line for monads on $\Cont$.

The new contribution (\S\ref{sec:enrichment}) is a clean anatomy of \emph{why} fullness holds or
fails, and a sharp collapse. There are three extensions of an $M$-container into endofunctors, with
three different notions of naturality:
\begin{itemize}[leftmargin=1.6em]
  \item[(A)] the \emph{ordinary} extension $\extM{-}\colon \Fam(\Kl(M)^{\op}) \to [\Set,\Set]$, whose
    morphisms are ordinary (pure) natural transformations;
  \item[(B)] the \emph{Kleisli-presheaf} extension $\extKl{-}\colon \Fam(\Kl(M)^{\op}) \to [\Kl(M),\Set]$,
    whose morphisms are Kleisli-natural, and which is \emph{always} full and faithful (co-Yoneda);
  \item[(C)] the \emph{self-enriched} extension $\ext{S,P}=\sum_s [P_s,-]\colon \C\to\C$ using the
    internal hom of a closed base, which realises Ghani's enriched-Yoneda desideratum
    $\Nat(\sum_s[P_s,-],F)\cong\prod_s F(P_s)$.
\end{itemize}
Two invariants govern (C), and they are independent:
\[
  \CL = \{\,M : \Kl(M)\text{ is monoidal closed}\,\}
  \qquad
  \UN = \{\,M : \text{the monoidal unit of }\Kl(M)\text{ is connected}\,\}.
\]
Closedness ($\CL$) makes the enriched extension \emph{exist}; connectedness of the monoidal unit
($\UN$) makes it \emph{full}. Our main theorem is that over a Kleisli base these two conditions meet in
exactly one monad.

\begin{theorem}[Enriched collapse; \S\ref{sec:enrichment}]\label{thm:main-intro}
Let $M$ be a commutative strong monad on $\Set$. The self-enriched $M$-container extension
$\sum_s[P_s,-]\colon\Kl(M)\to\Kl(M)$ is full and faithful if and only if $M$ is the identity monad.
Equivalently, $\CL\cap\UN=\{\Id\}$: no nontrivial effect --- not even an idempotent (localization)
monad --- carries a uniform enriched full-and-faithful semantics.
\end{theorem}

The reader monad $\Reader_S = (-)^S$ is the instructive pivot. Its Kleisli category is
\emph{cartesian closed}, so construction (C) is available; nonetheless the extension is not full,
because the monoidal unit $1$ has global-points functor $(-)^S$, which fails to preserve coproducts for
$|S|\ge 2$. The same failure appears in extension (A) as the non-codensity of $\Reader_S$. Closedness
and fullness are orthogonal; enrichment relocates the obstruction but does not remove it. The one
place fullness is automatic is the Kleisli-presheaf extension (B), at the price that its morphisms are
Kleisli-natural rather than pure. The factorisation $\extM{-}=R\circ\extKl{-}$, with $R$ restriction
along $F\colon\Set\to\Kl(M)$ and defect measured by the codensity monad $\Ran_M M$, is the single
picture carrying both messages.

\subsection*{Outline}
Section~\ref{sec:spine} sets up the change-of-base spine and fixes notation. Section~\ref{sec:census}
is the census of the six transport constructions, each stated with its provenance. Section
\ref{sec:enrichment} proves the anatomy and the collapse of Theorem~\ref{thm:main-intro}. Section
\ref{sec:conclusion} records what the diagnostic teaches, separates composition from the (lossy)
distributive-law question, and states the design consequence for effectful and agentic container
calculi.

%============================================================
\section{The change-of-base spine}\label{sec:spine}
%============================================================

We recall only what the later sections use.

\begin{definition}[Containers over a base]
Let $\C$ be a category. The category $\Fam(\C^{\op})$ of \emph{$\C$-containers} has objects
$(S,P)$ with $S\in\Set$ and $P\colon S\to\C_0$ a family of $\C$-objects $(P_s)_{s\in S}$; a morphism
$(S,P)\to(T,Q)$ is a function $f\colon S\to T$ on shapes together with, for each $s\in S$, a
$\C$-morphism $Q_{f(s)}\to P_s$ (contravariant in positions). For $\C=\Set$ this is $\Cont$.
\end{definition}

\begin{remark}[The bifibration viewpoint]\label{rmk:bifib}
The assignment $\C\mapsto\Fam(\C^{\op})$ is functorial: a functor $F\colon\C\to\D$ induces
$\Fam(F^{\op})\colon\Fam(\C^{\op})\to\Fam(\D^{\op})$, acting as $F$ on positions and identically on
shapes. The total category $\int_{\Set}\Cont(\mathrm{cod})$ assembling all bases is a bifibration over
$\Set$; contravariance in positions is exactly the fibrewise opposite. Change of base along
$F\colon\Set\to\C$ is the instance transporting the standard container theory into $\C$; when $F$ is
(co)monoidal it transports (co)monoidal structure, which is what the constructions of
\S\ref{sec:census} exploit.
\end{remark}

\begin{definition}[Extension]\label{def:extension}
When $\C$ is locally small the \emph{co-Yoneda (presheaf) extension} is
\[
  \extC{-}\colon \Fam(\C^{\op}) \longrightarrow [\C,\Set],\qquad
  \extC{S,P}(X) = \sum_{s\in S}\C(P_s, X),
\]
the coproduct of representables. When moreover a functor $M\colon\Set\to\Set$ makes
$\C(P_s,X)=\Set(P_s,MX)$ (as for $\C=\Kl(M)$), one also has the \emph{ordinary} extension
$\extM{-}\colon\Fam(\Kl(M)^{\op})\to[\Set,\Set]$, $\extM{S,P}(X)=\sum_s\Set(P_s,MX)$.
\end{definition}

\begin{remark}[Effectful reading]
For $\C=\Kl(M)$ the summand $\Set(P_s,MX)$ is the set of ways of assigning to each position an
$M$-effectful observation of $X$: $M=\Maybe$ makes positions optional, $M=\PP$ (powerset) makes them
nondeterministic, $M=\Dist$ probabilistic, $M=\Reader_S=(-)^S$ context-dependent. This is the sense in
which change of base to $\Kl(M)$ is ``the same datatype, read effectfully''.
\end{remark}

%============================================================
\section{A census of monad and comonad constructions}\label{sec:census}
%============================================================

We now collect the six ways monad/comonad structure has been transported through the container
construction. Each subsection states the result and its provenance; the proofs are in the cited
sources and are not reproduced. Throughout, $M$ is a monad on $\Set$ unless stated otherwise.

Two of the constructions (\S\ref{ssec:transfer}, \S\ref{ssec:liftings}) act \emph{fibrewise}: they
change the position family and leave the shape set alone, so they are transports of endostructure on
$\Set$ through the fibration $\Cont\to\Set$ of Remark~\ref{rmk:bifib}. Two more (\S\ref{ssec:mcont},
\S\ref{ssec:compose}) genuinely change the base or ask for structure \emph{on} $\Cont$. The fifth
(\S\ref{ssec:free}) is the free construction that turns a single container into a monad.

\subsection{Position-op transfer: monads become comonads}\label{ssec:transfer}
The simplest transport reads a monad on positions. Because positions are contravariant, a monad on
$\Set$ becomes a \emph{comonad} on $\Cont$.

\begin{construction}[Position transfer]\label{con:transfer}
For a monad $M=(M,\eta,\mu)$ on $\Set$, define $G\colon\Cont\to\Cont$ by
\[
  G(S,P)=(S,\,M\circ P),\qquad G(u,f)=(u,\,\{M f_s\}_s),
\]
with counit $\varepsilon_{(S,P)}=(\id_S,\,\eta_{P_s})$ (backward $P_s\to MP_s$) and comultiplication
$\delta_{(S,P)}=(\id_S,\,\mu_{P_s})$ (backward $MMP_s\to MP_s$).
\end{construction}

\begin{proposition}[\cite{MacBethCodensity}, transfer]\label{prop:transfer}
$(G,\varepsilon,\delta)$ is a comonad on $\Cont$, and this is a biconditional: the three comonad laws
for $G$ hold if and only if the three monad laws for $M$ hold (the single-shape containers $(1,A)$
realise every set $A$ as a fibre, so no law is lost). Dually, a comonad $W$ on $\Set$ yields a monad
$H(S,P)=(S,W\circ P)$ on $\Cont$. Fibrewise, $G$ is postcomposition by $M$ in the fibre
$(\Set^{\op})^S$; globally it is the left Kan extension
\[
  G(S,P) \;=\; \sum_{s\in S} \y^{\,M(P_s)} \;=\; \mathrm{Lan}_{(S,P)}\,M,
  \qquad
  \Cont\big(G(S,P),r\big)\;\cong\;[\Set,\Set]\big(\ext{S,P},\,r\circ M\big),
\]
naturally in $r$, and descends to $\ext{G(S,P)}(A)=\sum_s \Set(MP_s, A)$ on extensions.
\end{proposition}

The mechanism is exactly the slogan of the whole report in miniature: \emph{a monad on the base is a
comonad on the fibrewise opposite}, and $\Cont$ carries that opposite in its position variable
(Remark~\ref{rmk:bifib}).

\subsection{Predicate liftings: the $\forall/\exists$ constructions}\label{ssec:liftings}
A container $X=(S_X,P_X)$ and a ``predicate'' container $Y$ can be combined by quantifying the
predicate over positions. There are two ways, and they are the two halves of the composition monoidal
structure.

\begin{construction}[$\forall$ and $\exists$ liftings]\label{con:liftings}
For $(s,g)\in\ext{X}(S_Y)$ (so $s\in S_X$, $g\colon P_X s\to S_Y$), set
\[
  \mathrm{All}\,X\,Y\,(s,g)=\textstyle\prod_{p\in P_X s}P_Y(g\,p),
  \qquad
  \mathrm{Exists}\,X\,Y\,(s,g)=\textstyle\sum_{p\in P_X s}P_Y(g\,p),
\]
giving bifunctors $A\,X\,Y=(\ext{X}S_Y,\ \mathrm{All}\,X\,P_Y)$ and
$E\,X\,Y=(\ext{X}S_Y,\ \mathrm{Exists}\,X\,P_Y)$ on shapes $\sum_s S_Y^{P_X s}$.
\end{construction}

\begin{proposition}[\cite{MacBethCodensity}, $A/E$ liftings]\label{prop:liftings}
The $\exists$-lifting is exactly composition: $E = \tri$, a bifunctor on all of $\Cont$, with
$E\,X\,(E\,Y\,C)=E\,(X\tri Y)\,C$ the associativity of $\tri$ and unit $E\,\y\,C=C$. The
$\forall$-lifting $A$ is a bifunctor only on \emph{cartesian} morphisms in its first argument: a
natural pushforward along $\alpha=(u,\varphi)$ is indexed by set-sections of the backward map
$\varphi$, and the section-count is $0$, $\ge 2$, or exactly $1$ according as each $\varphi_s$ is
non-surjective, surjective-but-not-injective, or bijective. Hence $A$ extends canonically to
$\Cont_{\mathrm{cart}}\times\Cont\to\Cont$ only. On this domain $A$ is a left module of the
$\tri$-monoidal category $(\Cont,\tri,\y)$:
\[
  A\,X\,(A\,Y\,C)=A\,(X\tri Y)\,C=A\,(E\,X\,Y)\,C,\qquad A\,\y\,C=C.
\]
The mixed comparison $A\,X\,(E\,Y\,C)\to E\,(A\,X\,Y)\,C$ is an isomorphism iff $X$ is \emph{linear}
($|P_X s|=1$ for every $s$).
\end{proposition}

The cartesian restriction is precisely the phenomenon ``$A$ lifts through $\alpha$ iff $\alpha$ is
cartesian'', which for $\alpha=\mu_M$ is the classical criterion that $M$ lifts to containers iff $M$
is cartesian, seen one categorical level down. In the language of \S\ref{sec:spine}, $E=\Sigma$ is the
free (co-Yoneda) direction and $A=\Pi$ its right adjoint on the cartesian sublattice.

\begin{remark}[Lift / CoLift / BiLift]
Orestis' \textsc{Agda} development of the container effect library organises the same data as a
trichotomy \emph{Lift}, \emph{CoLift}, \emph{BiLift} of a monad through a container (Kodamai
\texttt{agda-containers}, \texttt{Containers/Effects}; personal communication). Proposition
\ref{prop:liftings} is the semantic content of that trichotomy: $E=\Sigma$ is the covariant (Lift)
direction, $A=\Pi$ the contravariant (CoLift) one, and their agreement is exactly the linear case.
\end{remark}

\subsection{$M$-containers as $\Fam(\Kl(M)^{\op})$, and codensity}\label{ssec:mcont}
The change of base to the Kleisli category is the central instance, and the one where the faithfulness
diagnostic first bites.

\begin{proposition}[\cite{MacBethCodensity}, Thm 1--2]\label{prop:mcont}
Let $M$ be a monad on $\Set$. There is an identity-on-data equivalence
$M\text{-}\Cont \simeq \Fam(\Kl(M)^{\op})$, and the effectful extension factors through the ordinary
one:
\[
  \extM{S,P} \;=\; \ext{S,P}\circ M,\qquad \extM{S,P}(X)=\sum_{s\in S}\Set(P_s, MX)
  =\sum_{s\in S}\Kl(M)(P_s,X),
\]
naturally in $(S,P)$. The extension $\extM{-}\colon M\text{-}\Cont\to[\Set,\Set]$ is \emph{always
faithful}, and
\[
  \extM{-}\text{ is full}\iff
  \underbrace{M\text{ is codense}}_{M\,\Rightarrow\,\Ran_M M\ \text{iso}}
  \ \wedge\
  \underbrace{\text{each }\Set(A,M-)\text{ is connected}.}_{\Ran_M\text{ preserves shape-coproducts}}
\]
If $M$ is affine ($M1\cong 1$) the connectedness clause holds automatically, so affine $+$ codense
$\Rightarrow$ full. Among polynomial effects $\Id$ is full while $\Maybe$, exceptions $X+E$, writer
$E\times X$ and reader $X^S$ ($|S|\ge2$) are not; the finite-distribution monad $\Dist$, being affine
and codense, \emph{is} full.
\end{proposition}

Two points deserve emphasis because they correct a natural but wrong guess. First, the criterion is
\emph{not} ``$M$ preserves coproducts'': that is the criterion for the coarser Kleisli-enriched target
of \S\ref{sec:enrichment}, and the right invariant for the target $[\Set,\Set]$ is codensity. Second,
the verdict is genuinely selective: probabilistic positions keep fullness, while error, reader, writer
and nondeterministic positions destroy it. The codensity monad $\Ran_M M$ is the exact measure of the
defect, a fact we use again in \S\ref{sec:enrichment}.

\subsection{The free monad on a container}\label{ssec:free}
Algebraic effects are presented by a signature, and a signature is a container: a set of operations
with an arity (a set of positions) each. The free monad on the signature functor is the type of
effect trees, and it is a monad \emph{on $\Cont$} built from a single container by iterated
composition. We record the linear instance, where the answer is cleanest and connects to a familiar
algebraic object.

\begin{proposition}[\cite{MacBethBase}, free $\tri$-monoid]\label{prop:free}
Let $C=(\star,A)$ be a one-shape linear container over a monoidal base $(\mathcal V,\otimes,k)$ (e.g.
$\mathcal V=\Vect$, $A$ finite-dimensional), so $\ext{C}(W)=A\otimes W$. The free $\tri$-monoid on $C$
is
\[
  C^{*}=\big(\mathbb N,\;(A^{\otimes n})_{n\in\mathbb N}\big),\qquad
  \ext{C^{*}}(W)=\bigoplus_{n\ge0}A^{\otimes n}\otimes W=T(A)\otimes W,
\]
with $T(A)$ the tensor algebra, and it carries the free-monad structure on $\ext{C}$: unit the
degree-$0$ injection and multiplication induced by concatenation $A^{\otimes m}\otimes A^{\otimes
n}\to A^{\otimes(m+n)}$. Thus the free monad on a linear container is the free monoid $T(A)$ acting by
$T(A)\otimes(-)$; over $\Vect$ this is exactly O'Neill's free monad, realised functorially.
\end{proposition}

\begin{remark}[The chain that computes it]
The free monad is the colimit of the Adámek chain of \emph{partial sums} $Y_n=\bigoplus_{k<n}
A^{\otimes k}\otimes W$, not of the bare powers $\ext{C}^{n}$: the endofunctor $A\otimes(-)$ is
pointable but never well-pointed, so the naive $\omega$-chain does not converge to the monad. The
pointing that repairs it is the two-shape container $(\{\mathrm{stop},\mathrm{layer}\},\
P_{\mathrm{stop}}=0,\ P_{\mathrm{layer}}=A)$ whose empty-position ``stop'' shape supplies the
degree-$0$ summand --- the residual/skip connection, read container-theoretically.
\end{remark}

More generally, the free monad on an arbitrary signature container is again a monad on $\Cont$
obtained by the same iterated-composition recipe: this is the standard reading of an algebraic
signature as a polynomial functor and its effect trees as the free monad, whose handler semantics has
a complete categorical account \cite{EffectHandlers}. Our $\tri$-monoid description is the
presentation-level shadow of that account.
% TODO(browse): deep-read Grodin, ``Poly-morphic effect handlers'' (Topos 2024), Neil's suggested
% reference for this item, then cite it here alongside \cite{EffectHandlers}. Not yet deep-read, so
% not cited per the provenance floor.

\subsection{When does a monad on $\Set$ give a monad on $\Cont$?}\label{ssec:compose}
The last construction is the subtlest, and it is where polynomiality --- \emph{representability} of
the effect --- enters. We ask when the essential image of $\extM{-}$ is closed under functor
composition, so that composing two effectful datatypes lands on another effectful datatype.

The obstruction is visible in one identity. For any $M$-containers $p,q$,
\[
  \extM{p}\circ\extM{q}
  \;=\;\big(\ext{p}\circ M\circ\ext{q}\big)\circ M
  \;=\; G\circ M,\qquad G:=\ext{p}\circ M\circ\ext{q},
\]
using the factorisation of Proposition~\ref{prop:mcont} twice. So \emph{every} composite of two
$M$-extensions already has the shape $G\circ M$; to land back in the essential image one must rewrite
the \emph{middle} $M\circ\ext{q}$ as (the extension of) a container. That is possible exactly when $M$
itself is a container --- i.e. when $M\cong\sum_{i\in I}\Set(B_i,-)$ is \emph{polynomial}
(representable) in the sense of \cite{GambinoKock}. This is the operational form of the composition question: \emph{closure under
composition is a representability condition on the middle effect}, no more and no less.

\begin{proposition}[\cite{MacBethCodensity}, composition]\label{prop:compose}
If $M$ is polynomial, with $\ceil{M}$ the container presenting it, then splicing $\ceil{M}$ in the
middle,
\[
  p\triM q \;:=\; p\tri\ceil{M}\tri q,\qquad
  \extM{\,p\triM q\,}\;\cong\;\extM{p}\circ\extM{q}\quad\text{naturally in }X,
\]
exhibits the essential image of $\extM{-}$ as closed under composition. The operation $\triM$ is
associative but has no unit up to container isomorphism unless $M=\Id$ (it is semigroupal, not
monoidal). Conversely, closure forces a cardinality growth law $|M(N(MX))|=R_N(|MX|)$ for a fixed
polynomial $R_N$; already at $p=q=\Id$ this excludes powerset $\PP$, nonempty powerset $\PP^{+}$ and
$\Dist$, none of which is polynomial.
\end{proposition}

The full converse --- closure $\Rightarrow$ $M$ polynomial --- is proved for the analytic monads of
non-flat symmetric operads in \cite{MacBethCodensity}, where the dividing line is a sharp
representability/rigidity condition on the middle effect; we do not reproduce that combinatorial
analysis here.

\begin{corollary}[Two independent invariants]\label{cor:indep}
Composability and fullness are logically independent. The writer monad $E\times(-)$ is polynomial (so
its $M$-containers compose) yet is not codense (so $\extM{-}$ is not full); the distribution monad
$\Dist$ is affine and codense (so $\extM{-}$ is full) yet is not polynomial (so its $M$-containers do
not compose). In particular the heuristic ``composes $\iff$ commutative and affine'' is false:
$\PP,\PP^{+},\Dist$ are commutative but fail composition, while non-affine $\Maybe$ and writer compose.
\end{corollary}

This is the first invariant of the effect: \emph{polynomiality governs composition}. The second,
\emph{codensity governs fullness}, is the theme of \S\ref{sec:enrichment}, where we show that passing
to an enriched base does not rescue it.

%============================================================
\section{Enrichment: three extensions and the collapse of fullness}\label{sec:enrichment}
%============================================================

The faithfulness diagnostic of \S\ref{ssec:mcont} was run against the target $[\Set,\Set]$. But
Ghani observed that when the base $\C$ is monoidal closed one can form a container extension valued in
$\C$ itself, using $\C$'s internal hom, and asked whether this \emph{enriched} extension enjoys a
uniform full-and-faithfulness theorem --- and, if so, whether it is available over the effectful bases
$\C=\Kl(M)$. The reader monad appears to make this urgent: its Kleisli category is cartesian closed
(so the enriched extension exists), yet \S\ref{ssec:mcont} said reader fails fullness. Is there a
contradiction?

There is not. The resolution is that \emph{three} distinct functors are in play, with three notions of
naturality; once separated, the reader is not full in any of them, consistently. We name them, prove
the availability criterion (Kock), tabulate the bases, and then prove the collapse.

\subsection{The three extensions}\label{ssec:three}
Fix a monad $M$ on $\Set$ and write $\C=\Kl(M)$.
\begin{itemize}[leftmargin=1.8em]
  \item[\textbf{(A)}] The \textbf{ordinary} extension
    $\extM{-}\colon \Fam(\Kl(M)^{\op})\to[\Set,\Set]$, $\extM{S,P}(X)=\sum_s\Set(P_s,MX)$. Morphisms of
    the target are \emph{pure} natural transformations. Faithful always; full iff $M$ codense and each
    $\Set(A,M-)$ connected (Proposition~\ref{prop:mcont}).
  \item[\textbf{(B)}] The \textbf{Kleisli-presheaf} extension
    $\extKl{-}\colon\Fam(\Kl(M)^{\op})\to[\Kl(M),\Set]$, $\extKl{S,P}=\sum_s\Kl(M)(P_s,-)$. This is the
    free-coproduct completion of $\Kl(M)^{\op}$: a coproduct of representables. Morphisms of the
    target are \emph{Kleisli}-natural. By ordinary co-Yoneda it is \emph{always} full and faithful,
    with no hypothesis on $M$.
  \item[\textbf{(C)}] Ghani's \textbf{self-enriched} extension. When $\C$ is symmetric monoidal
    closed with internal hom $[-,-]$, set
    \[
      \ext{S,P}=\sum_{s\in S}[P_s,-]\colon\C\longrightarrow\C,
    \]
    valued in $\C$-enriched endofunctors. Enriched co-Yoneda \cite{Kelly} gives Ghani's desideratum
    $\Nat\big(\sum_s[P_s,-],F\big)\cong\prod_s F(P_s)$ at the level of enriched hom-objects.
\end{itemize}
The three are related by a single factorisation, proved in \cite{MacBethCodensity}:
\[
  \extM{-}\;=\;R\circ\extKl{-},\qquad R=(-)\circ F\colon [\Kl(M),\Set]\to[\Set,\Set],
\]
where $F\colon\Set\to\Kl(M)$ is the free functor. On objects (A) and (B) agree; $R$ forgets
Kleisli-naturality, keeping only pure-naturality. \emph{All} loss of ordinary fullness lives in $R$,
and its size is measured exactly by the codensity monad $\Ran_M M$: pure-natural transformations
strictly contain the Kleisli-natural (equivalently, container) ones unless $M$ is codense.

\subsection{Availability: when is the base closed? (Kock)}\label{ssec:kock}
Extension (C) requires $\C=\Kl(M)$ to be closed, so we recall the criterion.

\begin{proposition}[Kock \cite{KockMonoidal,KockClosed}]\label{prop:kock}
Let $M$ be a strong monad on $\Set$.
\begin{enumerate}[label=\textup{(\arabic*)},leftmargin=2.4em]
  \item $\Kl(M)$ is symmetric monoidal (tensor $A\times B$ on objects, unit $1$, $F$ strict monoidal)
    \emph{iff} $M$ is commutative.
  \item $\Kl(M)$ is monoidal closed \emph{iff} the free algebras form an exponential ideal, i.e. each
    $(MB)^A\cong M[A,B]$ is free.
  \item $\Kl(M)$ is cartesian closed \emph{iff} additionally $M$ preserves finite products.
\end{enumerate}
Consequently a genuinely effectful monad whose Kleisli category has biproducts (hence a zero object)
is never cartesian closed --- a CCC with a zero object is trivial --- so the best such an effect can
reach is monoidal closure.
\end{proposition}

Write $\CL=\{M\ \text{commutative}: \Kl(M)\ \text{monoidal closed}\}$ for the availability class of
(C). Table~\ref{tab:bases} runs the change-of-base bases against these conditions; the ``monoidal
closed'' column is exactly where (C) is available.

\begin{table}[t]
\small
\begin{tabular}{@{}llll@{}}
\toprule
Base $\C$ & monoidal? & mon.\ closed? & cart.\ closed? \\
\midrule
$\Kl(\Maybe)=\Par$ & yes & yes & no \\
$\Kl(\Writer_E)$ & iff $E$ comm. & no (over $\Set$) & no \\
$\Kl(\Reader_S)=\mathrm{coKl}(-\times S)$ & yes & yes & \textbf{yes} \\
$\Kl(\PP)=\Rel$ & yes & yes & no \\
$\Kl(\Dist)=\Stoch$ & yes & no & no \\
$\Kl(\beta),\ \Kl(\mathrm{Cont}),\ \Kl(\List)$ & no & no & no \\
$\Vect$ \textit{(non-Kleisli)} & yes & yes & no \\
$\Cat,\ \Set$ \textit{(non-Kleisli)} & yes & yes & \textbf{yes} \\
\bottomrule
\end{tabular}
\medskip
\caption{Availability of the self-enriched extension across the change-of-base bases. ``Monoidal''
needs commutativity; $\beta$ (ultrafilter), continuation and list fail it, so are only premonoidal.
$\Stoch$ is monoidal but not closed ($(\Dist B)^A$ is a product of simplices, not a simplex).
$\Rel,\Vect$ are compact closed; $\Par$ is closed but not cartesian; only $\Set,\Cat,\Kl(\Reader)$ are
cartesian closed.}
\label{tab:bases}
\end{table}

\subsection{Fullness: connectedness of the unit}\label{ssec:un}
Availability is not fullness. By the flagship criterion of \cite{MacBethBase} (there written for a
general closed base), the self-enriched extension (C) is full and faithful iff the monoidal unit $I$
of $\C$ is \emph{connected}, i.e. $\C(I,-)$ preserves small coproducts --- this is what matches the
transformation set to the shape-coproduct $\sum_s$ in $\prod_s F(P_s)$. Write
$\UN=\{M: \text{the unit }1\text{ of }\Kl(M)\text{ is connected}\}$. Since $\Kl(M)(1,X)=\Set(1,MX)=MX$,
membership $M\in\UN$ says precisely that $M$ preserves small coproducts.

The two conditions are logically independent (Table~\ref{tab:2x2}), and this is the crux: closedness
and connectedness are orthogonal axes. Ghani's cartesian-closed reader sits in $\CL$ but not $\UN$;
the writer monad sits in $\UN$ but not $\CL$.

\begin{table}[t]
\small
\begin{tabular}{@{}lll@{}}
\toprule
 & connected ($\UN$) & disconnected \\
\midrule
closed ($\CL$) & $\Id$ \ \textbf{(target)} & $\Reader_S\,(|S|\ge2)$; $\Maybe$; $\PP$ \\
not closed & $\Writer_E$ ($E$ nontrivial) & $\Dist$; $\List$ \\
\bottomrule
\end{tabular}
\medskip
\caption{The two axes are independent. Each cell holds a commutative strong monad. The self-enriched
extension is available in the top row and full in the left column; only $\Id$ is in both.}
\label{tab:2x2}
\end{table}

\subsection{The collapse}\label{ssec:collapse}
We now prove that over a Kleisli base the two conditions meet only at the identity.

The engine of Step~2 below is the following elementary but base-independent fact, isolated here so that
no feature of $\Set$ is smuggled in.

\begin{lemma}[Closed $+$ initial $\Rightarrow$ terminal]\label{lem:closed-initial-terminal}
Let $(\mathcal{K},\otimes,I)$ be a symmetric monoidal closed category, with internal hom $[-,-]$
adjoint to the tensor via $\mathcal{K}(C\otimes A,B)\cong\mathcal{K}(C,[A,B])$ naturally, and suppose
$\mathcal{K}$ has an initial object $\emptyset$. Then $\mathcal{K}$ has a terminal object, namely
$[\emptyset,X]$ for \emph{any} object $X$ (canonically $\top:=[\emptyset,\emptyset]$).
\end{lemma}

\begin{proof}
For each object $A$ the functor $(-)\otimes A$ is a left adjoint (its right adjoint is $[A,-]$), hence
preserves all colimits, in particular the initial object; thus $C\otimes\emptyset\cong\emptyset$ for
every $C$. Fix any $X$. Naturally in $C$,
\[
  \mathcal{K}(C,[\emptyset,X])\;\cong\;\mathcal{K}(C\otimes\emptyset,X)\;\cong\;\mathcal{K}(\emptyset,X)\;\cong\;1,
\]
the last isomorphism because $\emptyset$ is initial. So $[\emptyset,X]$ is terminal. (In the symmetric
case the handedness of the closure is immaterial: $\emptyset\otimes C\cong\emptyset$ equally, so the
same computation goes through for a hom adjoint to left-tensoring.)
\end{proof}

Note that the lemma uses \emph{only} monoidal closedness together with an initial object: it invokes
neither cartesian closedness nor connectedness of the unit, which are the two independent axes of
\S\ref{ssec:un} (existence of the internal hom versus fullness of the extension).

\begin{theorem}[$\CL\cap\UN=\{\Id\}$]\label{thm:collapse}
Let $M$ be a commutative strong monad on $\Set$. The self-enriched $M$-container extension
$\sum_s[P_s,-]\colon\Kl(M)\to\Kl(M)$ is full and faithful if and only if $M\cong\Id$ as a monad.
\end{theorem}

\begin{proof}
By \S\ref{ssec:un}, fullness $\iff M\in\CL\cap\UN$. We show this intersection is $\{\Id\}$.

\emph{Step 1: $\UN$ forces the writer form.} Let $M\in\UN$, so $M$ preserves small coproducts. Put
$E:=M1$. For a set $T$, the assembly of the injections $M(\iota_t)\colon M1\to MT$ over $t\in T$
(where $\iota_t\colon 1\to T$ picks out $t$) is the copower-comparison map
\[
  \theta_T\colon T\times E\longrightarrow MT,\qquad \theta_T(t,e)=M(\iota_t)(e),
\]
which connectedness makes a bijection, naturally in $T$, with $\theta_1=\id_{M1}$. Hence
$M\cong(-)\times E$ as a functor. The monad structure on $(-)\times E$ is forced to be the writer
monad of a monoid structure $(E,m,e_0)$ (the monad laws are the monoid laws), commutative since $M$
is. Thus $\UN=\{\text{writer monads on commutative monoids}\}$, with $E=M1$; conversely every writer
monad preserves small coproducts because $\times$ distributes over $\coprod$ in $\Set$.

\emph{Step 2: $\CL$ forces $E\cong1$.} Let $M\in\CL\cap\UN$, so $MX\cong X\times E$. The base
$\Kl(M)$ is symmetric monoidal closed (Kock's tensor is the base product $\times$ with unit $1$, and it
is symmetric since $M$ is commutative) and has an initial object, namely the base initial object
$\emptyset$, because $\Kl(M)(\emptyset,B)=\Set(\emptyset,MB)=1$. Lemma~\ref{lem:closed-initial-terminal}
therefore supplies a terminal object $\top:=[\emptyset,\emptyset]$ in $\Kl(M)$. We stress that this step
uses \emph{only} monoidal closedness and the existence of an initial object; it uses neither cartesian
closedness nor connectedness of the unit --- closedness (existence) and connectedness (fullness) are the
two independent axes of \S\ref{ssec:un}. A terminal object
satisfies $\Set(C,M\top)=\Kl(M)(C,\top)\cong1$ for all $C$, forcing $M\top\cong1$. Using the writer
form, $\top\times E\cong M\top\cong1$, whence $|\top|\cdot|E|=1$ and $E\cong1$. Therefore
$MX\cong X\times1\cong X$, and the only monad structure on $\Id$ is the identity monad: $M\cong\Id$.

\emph{Step 3: $\Id$ qualifies.} $\Kl(\Id)=\Set$ is cartesian closed, so $\Id\in\CL$; and
$\Set(1,-)=\Id$ preserves all coproducts, so $\Id\in\UN$. Its self-enriched extension
$\sum_s(-)^{P_s}$ is the classical container extension, which is full and faithful. Hence
$\CL\cap\UN=\{\Id\}$.
\end{proof}

\begin{corollary}[No idempotent escape]\label{cor:idem}
No nontrivial idempotent (localization) monad on $\Set$ lies in $\CL\cap\UN$. Indeed by Step~1
membership in $\UN$ forces the writer form, and a writer monad is idempotent only when $E\cong1$; the
terminal monad $LX=1$ already fails connectedness, since $L(1+1)=1\ne 2=L1+L1$.
\end{corollary}

\begin{remark}[The two open gaps of the underlying criterion do not affect the collapse]
The flagship criterion of \cite{MacBethBase} proves necessity of unit-connectedness via the
\emph{copower} test, and whether that upgrades to preservation of \emph{all} small coproducts in every
closed base is left open there. It is irrelevant here: Step~1 uses only copower-preservation ($T\cdot
1$), which already forces the writer form, so every non-$\Id$ member of $\CL$ fails even the copower
test. Likewise, weakening $\UN$ to finite-coproduct preservation would still give $M1\cong1$ and
$Mn\cong n$ on finite sets; whether a commutative closed monad differing from $\Id$ only on infinite
sets could then exist is the one question a weaker definition would leave open, and we conjecture not.
Under the definition used here the classification is complete.
\end{remark}

\subsection{The reconciliation}\label{ssec:reconcile}
Theorem~\ref{thm:collapse} dissolves the apparent paradox. Ghani's claim is that the \emph{base}
$\Kl(\Reader_S)$ is cartesian closed --- so construction (C) is \emph{available}. It is not a claim
that any container extension is full. And in fact fullness fails in every picture, for one reason seen
through three lenses:
\begin{itemize}[leftmargin=1.8em]
  \item \textbf{(A) ordinary:} $\Reader_S$ is affine but not codense, so $\extM{-}$ is not full
    (Proposition~\ref{prop:mcont}).
  \item \textbf{(C) self-enriched over $\Kl(\Reader_S)$:} the base is closed, so (C) exists; but its
    unit $1$ has global-points functor $\Kl(\Reader_S)(1,-)=\Set(S,-)=(-)^S$, which fails to preserve
    coproducts for $|S|\ge2$. Disconnected unit $\Rightarrow$ (C) not full, by \S\ref{ssec:un}.
  \item \textbf{(B) Kleisli-presheaf:} always full and faithful --- but its morphisms are the
    Kleisli-natural ones, not the pure-natural transformations (A) asks about.
\end{itemize}
So cartesian closedness buys the \emph{availability} of an internal-hom extension, not its
\emph{fullness}. The reader's failure is enrichment-invariant: it is non-codensity in (A) and a
disconnected unit in (C), one obstruction wearing two masks. The dictionary is:
\[
  \begin{array}{lll}
    \textbf{extension} & \textbf{available} & \textbf{full and faithful} \iff\\[2pt]
    \text{(C) self-enriched }\sum_s[P_s,-] & \C\ \text{closed} & \text{unit of }\C\ \text{connected}\\
    \text{(B) Kleisli-presheaf }\extKl{-} & \text{always} & \text{always}\\
    \text{(A) ordinary }\extM{-} & \text{always} & M\ \text{codense and each }\Set(A,M-)\ \text{connected}
  \end{array}
\]
bridged by $\extM{-}=R\circ\extKl{-}$, whose defect is the codensity monad $\Ran_M M$ --- the exact
gap between the always-faithful enriched picture and the codensity-gated ordinary one.

%============================================================
\section{What the diagnostic teaches}\label{sec:conclusion}
%============================================================

Read as a whole, the census and the collapse say something about the change-of-base programme that no
single instance says on its own.

\emph{The frame is the content.} Six ostensibly different constructions --- position transfer,
$\forall/\exists$ liftings, Kleisli containers, free monads on signatures, and the composition of
effectful datatypes --- are one construction, $\Fam((-)^{\op})$, applied to different bases or read in
different variances. The position transfer (\S\ref{ssec:transfer}) and the $\forall/\exists$ liftings
(\S\ref{ssec:liftings}) are fibrewise transports; the Kleisli identification (\S\ref{ssec:mcont}) is
the genuine change of base; the free-monoid (\S\ref{ssec:free}) and composition
(\S\ref{ssec:compose}) results are about structure carried \emph{along} that change. Naming the frame
turns a list of theorems into a single diagram.

\emph{Full-and-faithfulness is a diagnostic, not a goal, and it exposes two invariants of an effect.}
Polynomiality (representability) governs \emph{composition}: the essential image of $\extM{-}$ is
closed under composition iff the middle effect is a container (\S\ref{ssec:compose}). Codensity governs
\emph{fullness} into $[\Set,\Set]$ (\S\ref{ssec:mcont}), and the same invariant reappears, under the
name unit-connectedness, as fullness of the enriched extension (\S\ref{ssec:un}). The two invariants
are independent (Corollary~\ref{cor:indep}): distributions compose no datatypes yet are full, writers
compose datatypes yet are not full.

\emph{Enrichment relocates the obstruction; it does not remove it.} This is the sharp lesson of
Theorem~\ref{thm:collapse}: passing to a closed base to gain Ghani's enriched-Yoneda formula does buy
the \emph{availability} of the self-enriched extension, but full-and-faithfulness then demands a
connected unit, and over a Kleisli base ``closed and connected'' pins the effect to the identity.
There is no effect monad over which one builds change-of-base containers $\Fam(\Kl(M)^{\op})$ with a
uniform full-and-faithful \emph{self-enriched} semantics. Any nontrivial effect must therefore be
carried either by the always-faithful Kleisli-presheaf extension (B) --- paying the price that its
transformations are Kleisli-natural rather than pure --- or by the ordinary $[\Set,\Set]$ extension
(A), with the codensity monad $\Ran_M M$ measuring the fullness defect. For a calculus of effectful or
agentic containers, where morphisms of containers are meant to be exactly the definable transformations
of the datatype, (B) is the honest target and codensity is the number to watch.

\subsection*{A different, lossier question}
Composition (\S\ref{ssec:compose}) asks when $\extM{p}\circ\extM{q}$ is again an $M$-extension. A
\emph{distinct} question --- when a distributive law $MM\Rightarrow M$ merges two layers of effect into
one --- is lossy: it collapses information that composition preserves, and its answer is governed by
different data. We flag the systematic comparison of the two as the natural next step, but it is not
the composition theorem of \S\ref{ssec:compose} in disguise.

\subsection*{Open questions and formalisation status}
\begin{enumerate}[label=(\arabic*),leftmargin=2.2em]
  \item The full converse of Proposition~\ref{prop:compose} (closure $\Rightarrow$ $M$ polynomial) is
    proved for the analytic monads of non-flat symmetric operads and open in general; the residual
    case is a representability question about the trailing effect, reduced to familial
    representability \cite{MacBethCodensity}.
  \item The two definitional gaps in the underlying fullness criterion (copower versus all
    coproducts; finite versus small coproduct connectedness) do not affect
    Theorem~\ref{thm:collapse}, but closing them would make the enriched criterion of
    \cite{MacBethBase} unconditional.
  \item The position transfer of \S\ref{ssec:transfer} is verified in Lean (the coordinate proof of
    the comonad laws, \texttt{[Quot.sound]}-only); the free $\tri$-monoid unit laws of
    \S\ref{ssec:free} are likewise machine-checked \cite{MacBethBase}. The collapse of
    Theorem~\ref{thm:collapse} is not yet formalised; its two elementary steps (writer form,
    terminal-object argument) are natural candidates.
\end{enumerate}

%============================================================
\begin{thebibliography}{99}

\bibitem{AAG} M.~Abbott, T.~Altenkirch, N.~Ghani,
\emph{Containers: constructing strictly positive types},
Theoret.\ Comput.\ Sci.\ \textbf{342} (2005), 3--27.

\bibitem{AhmanUustalu} D.~Ahman and T.~Uustalu,
\emph{Directed containers as categories},
Electron.\ Proc.\ Theor.\ Comput.\ Sci.\ \textbf{207} (2016), 89--98. arXiv:1604.01187.

\bibitem{EffectHandlers} S.~Kura,
\emph{On complete categorical semantics for effect handlers},
preprint, 2026. arXiv:2602.03275.

\bibitem{GambinoKock} N.~Gambino and J.~Kock,
\emph{Polynomial functors and polynomial monads},
Math.\ Proc.\ Cambridge Philos.\ Soc.\ \textbf{154} (2013), 153--192. arXiv:0906.4931.

\bibitem{Kelly} G.~M.~Kelly,
\emph{Basic Concepts of Enriched Category Theory},
London Math.\ Soc.\ Lecture Note Ser.\ \textbf{64}, Cambridge Univ.\ Press, 1982.

\bibitem{KockMonoidal} A.~Kock,
\emph{Monads on symmetric monoidal closed categories},
Arch.\ Math.\ (Basel) \textbf{21} (1970), 1--10.

\bibitem{KockClosed} A.~Kock,
\emph{Closed categories generated by commutative monads},
J.\ Austral.\ Math.\ Soc.\ \textbf{12} (1971), 405--424.

\bibitem{MacBethCodensity} MacBeth,
\emph{Two invariants of an effect monad: codensity governs fullness, polynomiality governs
composition}, Kodamai preprint, 2026.

\bibitem{MacBethBase} MacBeth,
\emph{Containers over a base}, Kodamai working paper, 2026.

\bibitem{NiuSpivak} N.~Niu and D.~I.~Spivak,
\emph{Polynomial Functors: A Mathematical Theory of Interaction},
book draft, 2023. arXiv:2312.00990.

\end{thebibliography}

\end{document}
